\section{Detection Workflows}
\label{sec:detection-strategies}

While the component architecture defines what responsibilities exist in the system, IDE usability depends on how these components are orchestrated in concrete workflows. Code Guardian supports five workflows that trade off latency, context breadth, and output surface. Table~\ref{tab:concept-workflows-overview} summarises them; the subsections that follow describe each in detail and relate them to requirements R1--R5. The process logic shared across the four detection workflows is illustrated by the BPMN in Section~\ref{subsec:process-view}; differences are confined to the trigger and the output surface.

\begin{table}[H]
  \centering
  \caption{Workflow modes supported by Code Guardian. The four detection workflows share the same backend pipeline (extract scope → optional RAG → Stage~1 LLM → parse → render); they differ only in trigger and output surface. Stage~2 repair is invoked per finding from any of the four detection workflows.}
  \label{tab:concept-workflows-overview}
  \footnotesize
  \renewcommand{\arraystretch}{1.2}
  \begin{tabularx}{\textwidth}{p{0.20\textwidth}p{0.20\textwidth}p{0.20\textwidth}X}
    \toprule
    \textbf{Workflow} & \textbf{Trigger} & \textbf{Scope} & \textbf{Output surface} \\
    \midrule
    Real-time function diagnostics (\S\ref{subsec:inline-detection}) & \texttt{onDidChangeText} (debounced 800\,ms) & Enclosing function ($\leq$2\,000 chars) & VS Code diagnostics, Problems panel, quick fixes \\
    On-demand file diagnostics (\S\ref{subsec:comprehensive-analysis}) & Command palette & Active document ($\leq$20\,000 chars) & VS Code diagnostics, Problems panel \\
    Interactive analysis \& Q\&A (\S\ref{subsec:interactive-analysis}) & Command palette, selection action & Selection / files \& folders & WebView panel (streamed Markdown) \\
    Workspace scan (\S\ref{subsec:workspace-scan}) & Command palette & \texttt{*.js/.ts/.tsx/.jsx} ($\leq$500\,KB / file) & Dashboard WebView, Problems panel \\
    \midrule
    Stage~2 repair (\S\ref{subsec:repair-workflow}) & Per Stage~1 finding & Single finding & VS Code quick fix (developer-applied) \\
    \bottomrule
  \end{tabularx}
\end{table}

\subsection{Real-Time Function-Level Diagnostics}
\label{subsec:inline-detection}

The real-time workflow provides continuous feedback while a developer edits JavaScript/TypeScript code. The key goal is to deliver timely, IDE-native warnings without disrupting flow (R4).

The extension listens to document change events for JavaScript and TypeScript documents, and analysis is \textbf{debounced} (default: 800\,ms, following common VS Code extension practice for balancing responsiveness against unnecessary re-analysis) so the model is not invoked on every keystroke. When the debounce fires, Code Guardian extracts the innermost enclosing function at the cursor position and analyzes only that snippet. A size guard skips unusually large functions (default: 2\,000 characters, chosen empirically to keep prompt size within the range where local models produce reliable structured output) to bound worst-case latency and avoid overloading the local model.

The analyzer is prompted to return a \textbf{strict JSON array of issues}, each containing a message and a line range. Repair suggestions are generated separately in Stage~2 and attached before rendering. The diagnostic adapter maps the snippet-relative, 1-based line numbers to VS Code ranges using the extractor's line offset and clamps ranges to valid document bounds. This enables stable rendering in the Problems panel, editor squiggles, and hover tooltips.

\textbf{Function-level analysis} improves responsiveness, but it can miss vulnerabilities whose evidence lies outside the current scope (e.g., validation in a different module). This limitation motivates the on-demand and workspace workflows and is addressed further in the future-work chapter.

\subsection{On-Demand File Diagnostics}
\label{subsec:comprehensive-analysis}

The file workflow provides broader coverage when a developer explicitly requests a deeper scan. A \textbf{command palette action} runs analysis over the full active document. To avoid generating excessively large prompts, the prototype skips very large files (default: 20{,}000 characters) and warns the user instead.

Findings are surfaced through the \textbf{same structured diagnostics pipeline} as the real-time workflow. This keeps the UI consistent, and it ensures that file scans can be reviewed in the Problems panel and navigated using standard IDE tooling.

\subsection{Interactive Analysis and Follow-up Q\&A}
\label{subsec:interactive-analysis}

Some security questions require richer explanations than a single diagnostic message. Code Guardian therefore includes webview-based interaction modes.

In the \textbf{selection analysis view}, the user can analyze a selected region (or the current line) and inspect the response in a dedicated analysis panel. This mode supports follow-up questions and streams Markdown-formatted responses for readability. Because it is user-initiated and not continuously triggered, it can tolerate higher latency than real-time diagnostics.

The \textbf{contextual Q\&A panel} allows the user to select files and folders as context and ask security questions about that subset of the workspace. The extension reads the selected files locally and sends their contents only to the local LLM backend, preserving the no-exfiltration objective for source code.

When enabled, the \textbf{RAG manager} can enrich prompts for these interactive workflows by injecting retrieved security knowledge snippets (CWE/OWASP/CVE-derived guidance). Retrieval and embeddings are performed locally; optional knowledge refresh operations fetch only public metadata.

\subsection{Workspace Scan and Security Dashboard}
\label{subsec:workspace-scan}

The workspace workflow supports periodic audits and prioritization across a repository.

The \textbf{file scanner} enumerates \texttt{*.js}, \texttt{*.jsx}, \texttt{*.ts}, and \texttt{*.tsx} files in the workspace while excluding dependency folders such as \texttt{node\_modules}. Very large files are skipped (default: $>$500\,KB) to keep runtime bounded.

Scan results are aggregated into a \textbf{dashboard WebView}. In addition to listing per-file issues, the dashboard computes a coarse security score based on issue density (issues per KLOC) and a keyword-based severity heuristic. This score is intended for prioritization rather than as a formal risk metric. The dashboard supports opening affected files directly from the report and rerunning scans, so the workflow complements real-time diagnostics by helping developers understand which parts of the codebase concentrate the most findings.

\subsection{Repair Suggestion Workflow}
\label{subsec:repair-workflow}

After consensus-filtered detection (Stage~1), the system generates targeted repairs for each confirmed vulnerability using a dedicated repair prompt (Stage~2). Repairs are exposed as VS Code quick fixes. Applying repairs is always user-initiated and integrates with the undo stack. This preserves developer control (R3) and reduces the risk of unintended behavioral changes from automatically applied patches.

\subsection{Confidence Abstention and Hybrid Gating}
\label{subsec:confidence-abstention}

Two design decisions complement consensus filtering to address the false-positive problem identified in the analysis chapter (R1) without re-introducing the deterministic-rule rigidity that motivated using LLMs in the first place.

\textbf{Multi-pass consensus} is itself an empirical signal. If two seeded passes report the same finding, the model agrees with itself; if only one pass reports it, the finding is noisier. The design promotes that signal from a binary keep-or-drop decision to an explicit per-finding confidence value, computable as the fraction of passes that produced a matching finding. A configurable threshold then decides which findings reach the diagnostic stream. The default behavior is no gating. Operators who care more about precision than recall can lift the threshold to suppress single-pass noise, which makes the precision/recall trade-off a configuration choice rather than a design constant.

A second source of corroboration is the established \textbf{SAST tooling}. When a deterministic baseline (Semgrep) flags a finding on overlapping lines and a matching canonical category, the LLM finding is promoted to maximum confidence and labeled as hybrid-source. The hybrid label is reported back to the IDE so the developer can see which findings have independent corroboration. This re-uses Semgrep's strengths (low false-positive rate on known patterns) without adopting its weaknesses (very narrow coverage), and it does so transparently --- the LLM's own findings remain visible at their original confidence even when Semgrep is silent.

Both mechanisms are conservative: they extend the existing pipeline with optional gates rather than replacing the un-gated path. The chapter's evaluation reports both with-gate and without-gate metrics so the contribution of each gate is attributable.

